\color{uniblue}\section[Сервисные функции]{Сервисные функции}\label{sec:servfuncs}
\color{black}

\needspace{4\baselineskip}
\color{uniblue}{\subsection{Группы уставок}}
\color{black}

\begin{enumerate}[label=\arabic{section}.\arabic{subsection}.\arabic*, labelsep=4pt, leftmargin=0pt, itemindent=57pt]
    
\item
В устройстве предусмотрено оперативное переключение между несколькими сохраненными группами уставок (ГУ).

\item
Оперативный вывод всех функций в составе устройства из работы происходит при появлении активного сигнала на входе <<Вывод терминала>>.
\item
Изменение группы уставок возможно:
\begin{itemize}
    \item по месту, через устройство ЮНИТ-ИЧМ;
    \item дистанционно, через команды от АСУ ТП;
    \item дистанционно, средствами сервисного ПО <<ЮНИТ-СЕРВИС>>;
    \item по месту, через соответствующий дискретный вход сигналом импульсного типа. Соответственно, использование ключа или кнопки с фиксацией не допускается.
\end{itemize}
\item
Переключение группы уставок происходит не мгновенно. После появления сигнала/команды на изменение активной группы уставок и до момента активации новой группы проходит от 7 до 15 с. В это
время происходит проверка целостности файлов выбранной ГУ. Всё это время устройство продолжает функционировать на той группе уставок, которая была активна до появления сигнала/команды на изменение активной группы уставок.

\end{enumerate}

%%%%%%%%%%%%%%%%%%%%%%%%%%%%%%%%%%%%%%%%%%%%%% АВАРИЙНЫЙ ОЦИЛЛОГРАФ %%%%%%%%%%%%%%%%%%%%%%%%%%%%%%%%%%%%%%%%%%%%%%%%%%
\needspace{4\baselineskip}
\color{uniblue}{\subsection{Аварийный осциллограф}}
\color{black}

\begin{enumerate}[label=\arabic{section}.\arabic{subsection}.\arabic*, labelsep=4pt, leftmargin=0pt, itemindent=57pt]

\item
При возникновении нарушений нормальных режимов работы данный процесс может записываться в память с помощью аварийного осциллографа. Осциллографирование выполняется с частотой дискретизации не менее 1200 Гц. Осциллограммы содержат измеренные значения аналоговых сигналов, дискретные сигналы состояния входов, выходных реле, а также логических сигналов функций РЗА и сервисные сигналы устройства.
\item
Пуск осциллографа производится по факту изменения состояния дискретных сигналов, заданных в качестве критериев пуска. В зависимости от длительности записи предаварийного и послеаварийного режима, в памяти может храниться разное количество осциллограмм. Данные, собранные аварийным осциллографом, сохраняются в энергонезависимой памяти устройства и циклически перезаписываются. Удаление записей аварийного осциллографа пользователем невозможно.

\end{enumerate}

%%%%%%%%%%%%%%%%%%%%%%%%%%%%%%%%%%%%%%%%%%%%%% ЖУРНАЛ СОБЫТИЙ %%%%%%%%%%%%%%%%%%%%%%%%%%%%%%%%%%%%%%%%%%%%%%%%%%
\needspace{4\baselineskip}
\color{uniblue}{\subsection{Журнал событий}}
\color{black}

\begin{enumerate}[label=\arabic{section}.\arabic{subsection}.\arabic*, labelsep=4pt, leftmargin=0pt, itemindent=57pt]

\item
устройстве реализована функция записи сигналов дискретных входов, работы защит, автоматики и сервисных сигналов в журнал событий.
\item
Максимальное число записей журнала - 20000. Каждая запись в журнал событий содержит в себе следующую информацию:
\begin{itemize}
    \item дату и время возникновения события;
    \item идентификатор события;
    \item состояние дискретного или значение аналогового сигнала в момент события;
    \item качество зарегистрированного параметра;
    \item качество времени;
\end{itemize}
\item
Данные о событиях сохраняются в энергонезависимой памяти. При заполнении журнала событий перезапись осуществляется от более старых событий к новым, то есть циклически. Удаление информации из журнала событий пользователем невозможно.

\end{enumerate}

%%%%%%%%%%%%%%%%%%%%%%%%%%%%%%%%%%%%%%%%%%%%%% РЕЖИМ ТЕСТИРОВАНИЯ %%%%%%%%%%%%%%%%%%%%%%%%%%%%%%%%%%%%%%%%%%%%%%%%%%
\needspace{4\baselineskip}
\color{uniblue}{\subsection{Режим тестирования}}
\color{black}

\begin{enumerate}[label=\arabic{section}.\arabic{subsection}.\arabic*, labelsep=4pt, leftmargin=0pt, itemindent=57pt]

\item
Режим тестирования предусмотрен для возможности оперативной проверки элементов/функций, таких как: ФСУ, измерительные органы, дискретные входы и выходные реле устройства.

\end{enumerate}